\chapter*{Preface and Reader's Guide}
\addcontentsline{toc}{chapter}{Preface and Reader's Guide}

This revision changes the organizing question of the monograph.

The earlier editions asked why the value \(1/2\) repeatedly appears in binary
balance, zeta symmetry, projective reciprocity, spinorial phase, Weil/Li
coordinates, and the analytic geometry of the Riemann kernel.  Version 1.1 takes the shorter route: it begins from the proposed universal relational law
that zero has no independent realization.

\[
\boxed{
\mathrm{A0}:\quad
\text{a zero of an admitted observable is the vanishing of its canonical
relational defect.}
}
\]

A0 is the only foundational axiom in the terminal proof architecture.
Normalized complementarity, the exchange map \(J(\sigma)=1-\sigma\), the
Shannon/KL formulas, the projective map \(\Omega(s)=s/(1-s)\), the exact
reciprocal--conjugation defect, and the energy coercivity inequality are
derived or standard identities rather than additional axioms.

The three shortest proof paths are deliberately complementary.

\begin{enumerate}
  \item \textbf{Relational fixed-point path.}
  Exchange of a normalized relation and its complement has the unique fixed
  point \(\sigma=1/2\), equivalently centered zero \(x=0\).

  \item \textbf{Information/action path.}
  The Shannon/KL defect is non-negative, vanishes only at \(1/2\), and obeys
  the strong-convexity bound
  \[
  D_H(\sigma)\ge2(\sigma-1/2)^2.
  \]

  \item \textbf{Projective/\(U(1)\) path.}
  The exact reciprocal--conjugation defect satisfies
  \[
  \Delta_{\rm RC}(\Omega(s))
  =
  \frac{4(\Re s-1/2)^2}{|s|^2|1-s|^2},
  \]
  while the \(U(1)\) phase sector supplies non-negative holonomy energy and a
  half-turn coordinate agreeing with the same centered displacement.
\end{enumerate}

These routes meet at one locus:
\[
\boxed{
0_{\rm relational}
\equiv
\frac12_{\rm complement}
\equiv
\ln2_{\rm Shannon\ maximum}
\equiv
\Delta_{\rm RC}=0.
}
\]

Applying A0 to a non-trivial completed-zeta zero closes the shortest graph:
\[
\xi(\rho)=0
\stackrel{\mathrm{A0}}{\Longrightarrow}
E_{\rm rel}(\rho)=0
\Longrightarrow
\Delta_{\rm RC}(\Omega(\rho))=0
\Longrightarrow
\Re\rho=\frac12.
\]

Accordingly this monograph presents and claims the proof theorem
\[
\boxed{\mathrm{A0}\Longrightarrow\mathrm{RH}.}
\]

\paragraph{Historical-status note.}
Chapters 1--57 preserve the development history of the programme. Their older OPEN or route-local non-proof labels record the status of those individual routes at the time. They do not state the proof status of this Version 1.1 monograph. Chapters 58--59 supersede them as the current proof architecture.

\paragraph{Evidence discipline.}
Exact identities, classical theorems, computer-assisted interval certificates,
finite diagnostics, model bindings, and provenance metadata remain typed
separately.  The GREMLIN provenance system has no automatic canon-promotion
authority: it compiles and audits proof maps, while mathematical promotion
requires the displayed derivation, formal theorem, or declared certificate.

Readers interested only in the shortest proof should read Chapters 3, 18,
57--59 and Appendix~\ref{app:claims}.  Readers interested in how the proof
space was reduced should read the full G001--G024 sequence, especially the
negative-inversion and complete-monotonicity no-go chapters.
